\documentclass[12pt]{article}
\usepackage[utf8]{inputenc}
\usepackage[T1]{fontenc}
\usepackage{hyperref}
\title{Shelta Language Resources for K-12 Curriculum Development}
\author{Compiled by Claude Sonnet 5.5 (claude-sonnet-5-5) for Andr\'as Kornai}
\date{}
\begin{document}
\maketitle

\section{Introduction}
Shelta (also called Cant, Gammon, Sheldru or Irish Traveller Cant) is a mixed language used by Irish Travellers. Its vocabulary derives largely from Irish, altered by reversal, consonant transposition, affixing and substitution. Its grammar is largely English. Sources describe it as a secret argot used mainly to hide meaning from outsiders, with a small lexicon (one source estimates about five hundred words). Speaker numbers are uncertain and the figures quoted are dated. The search found no school curriculum for Shelta, and a Mercator source states it is not taught in schools in the Republic of Ireland. The resources below are therefore mostly background material.

\section{Primary Educational Resources}
No primers, textbooks or classroom materials were found in search results. The only teaching-adjacent items found are the language cards and word lists from Friends, Families and Travellers:
\begin{itemize}
\item Irish Cant or De Gammon (Friends, Families and Travellers): word and phrase translations and language-card infographics. \url{https://www.gypsy-traveller.org/irish-cant-or-de-gammon/}
\end{itemize}

\section{Linguistic and Reference Materials}
\begin{itemize}
\item Shelta, Wikipedia overview. \url{https://en.wikipedia.org/wiki/Shelta}
\item Rosetta Project (Long Now) entry for Shelta (sth). \url{https://rosettapanglossia.longnow.org/index.php?title=Sth}
\item Academic article summary on Project MUSE. I did not verify its content beyond the search listing. \url{https://muse.jhu.edu/article/214634/summary}
\end{itemize}

\section{Religious and Community Resources}
\begin{itemize}
\item Friends, Families and Travellers, a community organisation that publishes the Cant material above. \url{https://www.gypsy-traveller.org/irish-cant-or-de-gammon/}
\item The Mercator wiki page on Shelta in Ireland mentions Traveller organisations (Irish Traveller Movement, Pavee Point). \url{https://wiki.mercator-research.eu/languages:shelta_in_ireland_under_construction_2}
\end{itemize}
No religious texts in Shelta were found.

\section{Digital and Audio Resources}
No audio, video or corpus resources were found in these searches.

\section{Curriculum Development Context}
Shelta is not taught in schools, and the Mercator source reports no teacher support. Its small lexicon and its use as a covert register mean that any curriculum would need to be developed with the Traveller community, who may regard the language as private. Heritage-oriented units could use the Friends, Families and Travellers word lists as a starting point. Little else exists in the sources found.

\end{document}